\documentclass[12pt]{article}
% Спільна преамбула для звітів (XeLaTeX)
\usepackage{fontspec}
\usepackage{polyglossia}
\setmainlanguage{ukrainian}
\setotherlanguage{english}
\setmainfont{DejaVuSerif}[Extension=.ttf,UprightFont=*,BoldFont=*-Bold,ItalicFont=*-Italic,BoldItalicFont=*-BoldItalic]
\setsansfont{DejaVuSans}[Extension=.ttf,UprightFont=*,BoldFont=*-Bold,ItalicFont=*-Oblique,BoldItalicFont=*-BoldOblique]
\setmonofont{DejaVuSansMono}[Extension=.ttf,UprightFont=*,BoldFont=*-Bold,ItalicFont=*-Oblique,BoldItalicFont=*-BoldOblique,Scale=0.9]
\usepackage[a4paper,margin=2.2cm]{geometry}
\usepackage{booktabs}
\usepackage{tabularx}
\usepackage{array}
\usepackage{enumitem}
\usepackage{xcolor}
\usepackage{tcolorbox}
\usepackage{graphicx}
\usepackage{setspace}
\usepackage[hidelinks]{hyperref}
\usepackage{titlesec}
\usepackage{fancyhdr}
\usepackage{multirow}
\usepackage{colortbl}
\setlength{\parskip}{4pt}
\setlength{\parindent}{0pt}
\onehalfspacing

\definecolor{accent}{HTML}{1F4E79}
\definecolor{sgreen}{HTML}{E2F0D9}
\definecolor{wred}{HTML}{FBE5D6}
\definecolor{oblue}{HTML}{DEEBF7}
\definecolor{tyellow}{HTML}{FFF2CC}
\definecolor{lightgray}{HTML}{F2F2F2}

\titleformat{\section}{\large\bfseries\color{accent}}{\thesection}{0.6em}{}
\titleformat{\subsection}{\normalsize\bfseries}{\thesubsection}{0.6em}{}

\pagestyle{fancy}
\fancyhf{}
\lhead{\small ІТ, Модуль 2 — командний проєкт}
\rhead{\small Петров Б., Мекшаков А.}
\cfoot{\thepage}
\renewcommand{\headrulewidth}{0.3pt}

\newcolumntype{L}[1]{>{\raggedright\arraybackslash}p{#1}}
\newcolumntype{Y}{>{\raggedright\arraybackslash}X}

% Титульна сторінка
\newcommand{\maketitlepage}[3]{%
\begin{titlepage}
\centering
{\large Київський національний університет імені Тараса Шевченка\\[2pt]
Факультет комп'ютерних наук та кібернетики\par}
\vspace{3cm}
{\large Дисципліна: \textbf{Інформаційні технології}\par}
\vspace{0.5cm}
{\large Модуль 2 — робота в команді\par}
\vspace{2cm}
{\LARGE\bfseries #1\par}
\vspace{0.6cm}
{\large #2\par}
\vfill
\begin{flushright}
\begin{tabular}{l}
\textbf{Виконали:} студенти 4 курсу, група МІ-41\\[4pt]
Петров Богдан\\
Мекшаков Артемій\\[10pt]
\textbf{Спеціальність:} Комп'ютерні науки\\
\end{tabular}
\end{flushright}
\vspace{1.5cm}
{Київ — #3\par}
\end{titlepage}
}


\begin{document}
\maketitlepage{Звіт за етапом 0}{Формування команди: склад і розподіл ролей}{2026}

\section{Мета етапу}
Метою нульового етапу є формування команди для виконання командного проєкту в межах Модуля~2 дисципліни «Інформаційні технології», попередній розподіл ролей та узгодження правил спільної роботи. Склад команди фіксується на лабораторному занятті 21.10.2026.

\section{Склад команди}
Команда складається з двох учасників, що відповідає вимозі щодо чисельності (до 5--6 осіб включно). Невеликий розмір команди обрано свідомо: він спрощує комунікацію, усуває необхідність складної координації та дозволяє кожному учаснику охопити повний цикл розробки у своїй зоні відповідальності.

\begin{table}[h]
\centering
\renewcommand{\arraystretch}{1.3}
\begin{tabularx}{\textwidth}{L{0.6cm} L{4.2cm} Y}
\toprule
\textbf{№} & \textbf{Учасник} & \textbf{Основна роль (попередньо)}\\
\midrule
1 & Петров Богдан & Team Lead, Backend-розробник (NestJS), відповідальний за архітектуру та базу даних\\
2 & Мекшаков Артемій & Product Owner, Frontend-розробник (Next.js), відповідальний за UI/UX та якість (QA)\\
\bottomrule
\end{tabularx}
\caption{Склад команди та основні ролі}
\end{table}

\section{Розподіл ролей та зон відповідальності}
Оскільки команда складається з двох осіб, кожен учасник поєднує кілька ролей. Розподіл побудовано так, щоб навантаження було рівномірним, а зони відповідальності — чітко розмежованими за шарами системи (серверна частина / клієнтська частина) та за управлінськими функціями (технічне лідерство / продуктове бачення).

\begin{table}[h]
\centering
\renewcommand{\arraystretch}{1.3}
\begin{tabularx}{\textwidth}{L{3.6cm} Y Y}
\toprule
\textbf{Напрям} & \textbf{Петров Богдан} & \textbf{Мекшаков Артемій}\\
\midrule
Управління & Team Lead: технічні рішення, код-рев'ю, контроль дедлайнів & Product Owner: вимоги, пріоритезація беклогу, ведення дошки задач\\
\addlinespace
Розробка & Backend: NestJS, REST API, бізнес-логіка, автентифікація & Frontend: Next.js, компоненти інтерфейсу, інтеграція з API\\
\addlinespace
Дані та інфраструктура & Проєктування схеми БД, міграції, Docker, CI/CD & Макети інтерфейсу, дизайн-система, адаптивна верстка\\
\addlinespace
Якість & Модульні тести backend, API-документація (Swagger) & Інтеграційні/E2E-тести, тести UI, журнал тестування\\
\addlinespace
Документація та звітність & UML-моделювання, обґрунтування архітектури & Презентації, Lean Canvas, карта продукту, звіти за етапами\\
\bottomrule
\end{tabularx}
\caption{Зони відповідальності учасників}
\end{table}

\section{Правила спільної роботи}
\begin{itemize}[leftmargin=1.5em]
  \item \textbf{Система контролю версій:} GitHub, один репозиторій; робота через гілки \texttt{feature/*} та pull request'и з обов'язковим рев'ю другого учасника.
  \item \textbf{Управління задачами:} дошка в одній з платформ (Trello / Jira), спринти тривалістю один тиждень відповідно до графіка лабораторних занять.
  \item \textbf{Комунікація:} щоденні короткі синхронізації в месенджері, щотижневе демо та ретроспектива перед лабораторним заняттям.
  \item \textbf{Технологічний стек (попередньо):} TypeScript, Next.js (frontend), NestJS (backend), PostgreSQL, Docker.
  \item \textbf{Прийняття рішень:} технічні рішення — за Team Lead, продуктові — за Product Owner; спірні питання вирішуються спільно з урахуванням зворотного зв'язку викладача.
\end{itemize}

\section{Результат етапу}
Команду сформовано, попередні ролі розподілено, правила спільної роботи узгоджено. Подальші кроки: підготовка варіантів теми проєкту та їх SWOT-аналіз (етап~1), фіксація складу команди на лабораторному занятті 21.10.2026.

\end{document}
